\Needspace{12\baselineskip}
\begin{center}
\textit{Forms with square determinant.}
\end{center}

\Needspace{8\baselineskip}
\Article{206}
\textsc{Problem.} \textit{Given a form $(a,b,c)$ with square determinant $hh$, where $h$ denotes its positive root, find a properly equivalent form $(A,B,C)$ in which $A$ lies between $0$ and $2h-1$, inclusive, $B$ $=h$, and $C=0$.}

\textit{Solution.} I. Since $hh=bb-ac$, we have $(h-b):a=c:-(h+b)$. Let $\beta:\delta$ be an equal ratio, with $\beta$ relatively prime to $\delta$, and determine $\alpha$, $\gamma$ so that $\alpha\delta-\beta\gamma=1$, which is possible. Let the substitution $\alpha,\beta,\gamma,\delta$ take $(a,b,c)$ into $(a',b',c')$, which is therefore properly equivalent to it. We have
\[
\begin{aligned}
b'&=a\alpha\beta+b(\alpha\delta+\beta\gamma)+c\gamma\delta\\
&=(h-b)\alpha\delta+b(\alpha\delta+\beta\gamma)-(h+b)\beta\gamma\\
&=h(\alpha\delta-\beta\gamma)=h\\
c'&=a\beta\beta+2b\beta\delta+c\delta\delta\\
&=(h-b)\beta\delta+2b\beta\delta-(h+b)\beta\delta=0
\end{aligned}
\]
So if $a'$ also already lies between $0$ and $2h-1$, the form $(a',b',c')$ satisfies all the conditions.

II. But if $a'$ lies outside the bounds $0$ and $2h-1$, let $A$ be the least positive residue of $a'$ modulo $2h$, which clearly lies between these bounds, and put $A-a'=2hk$. Then the form $(a',b',c')$, \textit{i.e.}, $(a',h,0)$, passes by the substitution $1,0,k,1$ into $(A,h,0)$, which is properly equivalent to $(a',b',c')$, $(a,b,c)$ and satisfies all the conditions. --- Clearly the form $(a,b,c)$ passes into $(A,h,0)$ by the substitution $\alpha+\beta k,\beta,\gamma+\delta k,\delta$.

\textit{Example.} Let the given form be $(27,15,8)$, with determinant $=9$. Here $h=3$; the ratio in smallest terms equal to $-12:27=8:-18$ is $4:-9$. So putting $\beta=4$, $\delta=-9$, $\alpha=-1$, $\gamma=2$, the form $(a',b',c')$ becomes $(-1,3,0)$, which passes into $(5,3,0)$ by the substitution $1,0,1,1$. This is therefore
\SourceLeaf{206}
the required form, and the given form passes into it by the proper substitution $3,4,-7,-9$.

We will call such forms $(A,B,C)$, in which $C=0$, $B=h$, and $A$ lies between $0$ and $2h-1$, \textit{reduced forms}. They must therefore be carefully distinguished from reduced forms with negative determinant or positive nonsquare determinant.

\Needspace{8\baselineskip}
\Article{207}
\textsc{Theorem.} \textit{Two reduced forms $(a,h,0)$, $(a',h,0)$ that are not identical cannot be properly equivalent.}

\textit{Proof.} Suppose they are properly equivalent, and let the former pass into the latter by the proper substitution $\alpha,\beta,\gamma,\delta$. We then have four equations:
\begin{align*}
a\alpha\alpha+2h\alpha\gamma&=a'\tag*{. . . . . . . . . [1]}\\
a\alpha\beta+h(\alpha\delta+\beta\gamma)&=h\tag*{. . . . . . . [2]}\\
a\beta\beta+2h\beta\delta&=0\tag*{. . . . . . . . . . [3]}\\
\alpha\delta-\beta\gamma&=1\tag*{. . . . . . . . . . . [4]}
\end{align*}
Multiplying the second equation by $\beta$, the third by $\alpha$, and subtracting gives $-h(\alpha\delta-\beta\gamma)\beta=\beta h$, or, by [4], $-\beta h=\beta h$, so we must have $\beta=0$. Thus [4] gives $\alpha\delta=1$ and $\alpha=\pm1$. Then [1] gives $a\pm2\gamma h=a'$. This equation cannot hold unless $\gamma=0$ (since by hypothesis both $a$ and $a'$ lie between $0$ and $2h-1$), \textit{i.e.}, unless $a=a'$, or the forms $(a,h,0)$, $(a',h,0)$ are identical, contrary to the hypothesis.

We can therefore easily solve the following problems, which caused much greater difficulty for nonsquare determinants.

I. \textit{Given two forms $F$, $F'$ with the same square determinant, determine whether they are properly equivalent.} Find two reduced forms properly equivalent to $F$, $F'$ respectively. If they are identical, the given forms are properly equivalent; otherwise they are not.

II. \textit{Under the same assumptions, determine whether they are improperly equivalent.} Let the form opposite to one of the given forms, \textit{e.g.}, $F$, be $G$. If it is properly equivalent to $F'$, then $F$ and $F'$ are improperly equivalent, and conversely.
